\section{NN-A001: exact arithmetic label and a certified correction}
\label{sec:NN-A001}

\paragraph{Inputs.}
A supplies RH.  The analytic inputs are the Riemann--von Mangoldt identity
$N(t)=\theta(t)/\pi+1+S(t)$ with midpoint convention at a zero, and the
Hurwitz derivative identity DLMF 25.11.17.  The supplied consult
\path{OURS_CONSULT_ad855737_ZETA_LABELS.md}, Q1, Proposition 1, already
displays the correction term in its proof, but omits it from its claimed
equivalence.  We retain that proof's term and independently certify a
counterexample to the simplified assertion.

\begin{theorem}[NN-A001]
Let $\rho_n=1/2+i\gamma_n$ be a simple positive-ordinate zero, indexed with
multiplicity among all positive-ordinate zeros, and set
\[
 Z(t)=e^{i\theta(t)}\zeta(1/2+it),\quad
 \phi_n=\arg\zeta(3/2+i\gamma_n),\quad
 W_n=S(\gamma_n)-\phi_n/\pi.
\]
Choose the argument on the absolutely convergent Euler-product branch.
For the deformation $\zeta(s,1+t)$ put
$v_n=\rho_n\zeta(\rho_n+1)/\zeta'(\rho_n)$ and
$\delta_n=\arctan(1/(2\gamma_n))$.  Then
\[
 \operatorname{Re}v_n=
 \frac{|\zeta(\rho_n+1)|}{|Z'(\gamma_n)|}
 \left(\gamma_n\sin\pi W_n+\frac12\cos\pi W_n\right).       \tag{A001.1}
\]
Thus the exact right-mover criterion is
$\sin(\pi W_n+\delta_n)>0$, equivalently
$W_n+\delta_n/\pi\bmod2\in(0,1)$.  Tangency is exactly equality to an
integer modulo one after this correction.  The uncorrected criterion in
the prompt is false: the indexed zero $n=1169$ is a certified right-mover
although $W_{1169}<0$ and $\sin(\pi W_{1169})<0$.
\end{theorem}

\begin{proof}
Implicit differentiation using DLMF 25.11.17 gives the stated $v_n$.
Differentiating $Z(t)$ at a zero gives
$\zeta'(\rho_n)=-i e^{-i\theta(\gamma_n)}Z'(\gamma_n)$.
Since $Z(0)<0$, the sign of $Z'$ at this simple zero is
$(-1)^{n+1}$; preceding zeros of even multiplicity change neither this
sign nor the parity of the total zero count.  The midpoint convention
gives $\theta(\gamma_n)/\pi=n-3/2-S(\gamma_n)$.  Hence
\[
 \cos(\theta+\phi)=(-1)^n\sin\pi W_n,\qquad
 \sin(\theta+\phi)=(-1)^n\cos\pi W_n.
\]
Taking the real part of
$i\rho_n e^{i\theta}\zeta(\rho_n+1)/Z'(\gamma_n)$ proves (A001.1).
The positive factor $\sqrt{\gamma_n^2+1/4}$ converts its last parenthesis
to $\sqrt{\gamma_n^2+1/4}\sin(\pi W_n+\delta_n)$.

The included program \path{check_label_1169.py} calls the indexed
FLINT/Arb zero isolator and evaluates the derivative through a complex
ball power series, at 60 decimal digits.  It certifies
\[
 \gamma_{1169}=1613.0788554510274235\ldots,\quad
 W_{1169}=-0.0000606299724187599509\ldots,
\]
\[
 \delta_{1169}/\pi=0.0000986653178528222938\ldots,
 \qquad \operatorname{Re}v_{1169}=0.025925651967163814\ldots>0.
\]
The saved ball intervals in \path{label_certificate.json} exclude zero
for both sine signs and for the velocity's real part.  The program asserts
the relevant inequalities using ball comparisons, rather than midpoint
rounding.  This supplies the counterexample claimed in the theorem.
\end{proof}

\paragraph{Interpretation and geometry.}
The formula is a real normal component of a tangent vector to a moving
analytic Cartier divisor.  Its sign is a Boolean-valued statistic where
the normal component is nonzero; it is not a homomorphism of the compact
circle into a two-element group.  Indeed $\operatorname{sgn}\sin\pi x$ is
a discontinuous square wave, and even as a set map is not multiplicative
under addition: two small positive $x,y$ give $+,+,+$, while suitable
$x,y$ give $+,+,-$.  In particular the consult's description as a sign
character of a quotient circle is not a representation-theoretic identity.
The weight-zero Boolean skeleton in the split-zero construction contains
no value of $\zeta'(\rho)$ or of this normal vector.  This corrected label
therefore supplies no SAT, parity, or multiplicity-positivity algorithm.

\section{NN-A002: the exact statistical gate and why a CLT does not close it}
\label{sec:NN-A002}

\begin{theorem}[NN-A002]
A central limit theorem for $S(\gamma_n)$ after division by a standard
deviation tending to infinity, even together with boundedness of
$\phi_n$, does not imply that $W_n\bmod2$ equidistributes, nor that half
the labels in NN-A001 are positive.

A sufficient arithmetic gate for the half-proportion, on the chosen
simple-zero subsequence, is the following pair of assertions.  For every
odd integer $k$, the empirical means of $e^{i\pi kW_n}$ tend to zero;
and
\[
 \lim_{\epsilon\downarrow0}\limsup_{T\to\infty}
 \frac{\#\{n:\gamma_n\leq T,\ n\text{ simple},\
                 \operatorname{dist}(W_n,\mathbb Z)<\epsilon\}}
      {\#\{n:\gamma_n\leq T,\ n\text{ simple}\}}=0.       \tag{A002.1}
\]
All nonzero Fourier modes tending to zero is the stronger usual Weyl
equidistribution gate.  The required odd modes have the exact form
\[
 e^{i\pi kW_n}=e^{-3\pi i k/2}(-1)^n
          e^{-ik\theta(\gamma_n)}e^{-ik\phi_n}\qquad(k\text{ odd}).
                                                               \tag{A002.2}
\]
\end{theorem}

\begin{proof}
Let $G$ be standard normal, let $\sigma_j\to\infty$, and set
$X_j=2\lfloor\sigma_jG/2\rfloor+1/4$.
Then $|X_j-\sigma_jG|\leq7/4$, so $X_j/\sigma_j$ converges almost surely
and in every fixed $L^p$ to $G$.  In particular the normalized central
limit law and all its fixed normalized moment consequences hold.
Nevertheless $X_j\bmod2=1/4$ identically, so every square-wave label is
positive.  Equivalently, an exactly Gaussian variable $\sigma_jG$ can
be altered by a uniformly bounded, correlated correction to have this
completely nonuniform residue class.  This proves the failure of the
implication invoked in the supplied consult, Q2, Proposition 4.

For the sufficient gate, approximate the square wave
$f(x)=\operatorname{sgn}\sin\pi x$ by its Fejer trigonometric means.
It has only odd Fourier modes, because $f(x+1)=-f(x)$, and zero constant
term.  On any closed set a positive distance from $\mathbb Z$, these
means converge uniformly to $f$, and they are bounded in absolute value
by one.  For fixed truncation the empirical mean tends to zero by the
odd-mode hypothesis.  Equation (A002.1) controls the omitted boundary
neighborhood; then let the truncation tend to infinity and that
neighborhood shrink.  Thus the empirical mean of $f(W_n)$ tends to zero.
Since $\delta_n\to0$, (A002.1) also shows that replacing $f(W_n)$ by
$f(W_n+\delta_n/\pi)$ changes an asymptotically negligible proportion.
This proves the half-proportion for the exact labels.  Finally substitute
$W_n=n-3/2-\theta(\gamma_n)/\pi-\phi_n/\pi$ to get (A002.2).
\end{proof}

\paragraph{A primary CLT locator and its frequency range.}
The previously missing source can be made precise: Korolev,
\emph{On the Gram's Law in the Theory of Riemann Zeta Function},
arXiv:1011.3997, Theorem~14, printed pp.33--34, gives a Gaussian
law for $S(\gamma_n\pm0)/\sigma_N$, where
$\sigma_N=\sqrt{L/2}/\pi$ and $L=\log\log N$, on the stated
index blocks. Its proof invokes Theorem~10, pp.21--22. There the
characteristic-function argument is truncated at normalized frequency
$\lambda\asymp\sqrt{\log L}$. In the original $S$ coordinate
this reaches only $\lambda/\sigma_N
=O(\sqrt{\log L/L})\to0$, not the fixed nonzero frequencies
$\pi k$ required by NN-A002. The theorem also does not insert the
correlated exterior phase. Thus the concrete classical CLT, and not
merely an abstract CLT, leaves the identified fixed-frequency gate open.

\paragraph{Located progress beyond the invalid CLT route.}
Arias de Reyna, \emph{On the distribution (mod 1) of the normalized zeros
of the Riemann zeta-function}, arXiv:1403.7126, Theorem 3.1 and Corollary
3.2, prove cancellation of
$N(T)^{-1}\sum_{0<\gamma\leq T}e^{2i\kappa\theta(\gamma)}$ for
$0<\kappa<6/5$ once A identifies the ordinates with the real arguments
in that theorem.  The theorem itself is formulated without RH using
complex ordinates.  Its statement does not supply the alternating index
factor and exterior phase in (A002.2), does not select the simple-zero
subsequence, and does not establish all the required modes.  These are
specific additional gates.  The weighted first moment of velocities is
also insufficient: cancellation of a sum of real velocities controls
their sizes with signs, whereas the question here counts signs without
those weights.

\paragraph{Dictionary and complexity impact.}
The Fourier gate is a statement about characters $e^{i\pi kx}$ of a
compact real torus and a discontinuous observable on it.  The sample
points are analytic divisor points; no finite-type scheme or Frobenius
action on these phase values has been constructed.  This statement
neither establishes a nonstandard Weil purity theorem nor changes the
complexity of support, parity, or satisfiability.
